\section{Complete Dynamic Epistemic Reachability}
\label{sec:reachability}

The strongest endpoint theorem uses a uniform six-world construction scheme. For arbitrary source and target values
\[
  s,t\in\{\Tval,\Fval,\Bval,\Nval\},
\]
consider worlds
\[
  \{f,a_1,a_2,a_3,a_4,r\}.
\]
The distinguished focus world \(f\) carries the target value \(t\) for atom \(p\), each of the four anchor worlds \(a_k\) carries the source value \(s\), and the spare world \(r\) carries \(\Nval\). A second atom \(e\) is \(\Tval\) only at \(f\) and \(\Fval\) elsewhere. Every world accesses all six worlds. The prior measure is uniform,
\[
  \mu(\{w\})=\frac16,
\]
and the Lockean threshold is
\[
  c=\frac23.
\]

Before update, the four anchors contribute total mass \(4/6=2/3\). Coordinate by coordinate, they therefore force the threshold belief \(\Bel p\) to reproduce the source value \(s\), regardless of the single target-bearing focus world. After conditioning on \(e\), all posterior mass is concentrated at \(f\), so \(\Bel p\) reproduces the target value \(t\).

The local probability function is independent of the evaluation world. Consequently, the complete \(\Bel p\) profile is homogeneous across the common accessible range before and after update. The outer knowledge operator therefore passes both values through unchanged.

\begin{theorem}[Complete dynamic epistemic reachability]
For every ordered pair \((s,t)\) of complete FDE values, there is a member of the six-world family and an admissible conditionalization such that
\[
  \Know(\Bel p):s\longrightarrow t.
\]
Across that update the accessibility relation, atomic valuation, and threshold are fixed; only the local probability field changes.
\statuswitness
\end{theorem}

\begin{proof}[Proof sketch]
For either sign, a true source bit contributes at least \(4/6=c\), while a false source bit receives at most \(1/6<c\) from the focus. The evidence has mass \(1/6\); conditioning gives focus mass one. Homogeneity of the inner belief value makes each of its outer signed support events either the whole range or empty, so \(\Know(\Bel p)\) has the asserted endpoints. The source bits equal to the threshold use the inclusive \(\ge\) convention.
\end{proof}

Equivalently, the categorical reachability graph is complete:

\begin{center}
\begin{tabular}{c|cccc}
\toprule
source \textbackslash\ target & \(\Tval\) & \(\Fval\) & \(\Bval\) & \(\Nval\)\\
\midrule
\(\Tval\) & $\checkmark$ & $\checkmark$ & $\checkmark$ & $\checkmark$\\
\(\Fval\) & $\checkmark$ & $\checkmark$ & $\checkmark$ & $\checkmark$\\
\(\Bval\) & $\checkmark$ & $\checkmark$ & $\checkmark$ & $\checkmark$\\
\(\Nval\) & $\checkmark$ & $\checkmark$ & $\checkmark$ & $\checkmark$\\
\bottomrule
\end{tabular}
\end{center}

The theorem immediately gives creation and removal of both nonclassical statuses. In particular, the development contains verified instances
\[
  \Tval\to\Bval,\qquad
  \Bval\to\Nval,\qquad
  \Nval\to\Tval.
\]
Thus neither glutty nor gappy knowledge is dynamically terminal under the present update semantics.

\paragraph{Quantifier discipline.}
The theorem does not say that one fixed concrete model realizes all sixteen arrows. The family is parameterized by the requested source and target values. For each pair, however, the pre-update and post-update models share the same \(R\), \(v\), and \(c\), and the transition is generated by admissible probability-only conditionalization. The result is therefore a semantic non-prohibition theorem: no ordered pair of complete four-valued knowledge statuses is ruled out merely by the conditionalization mechanism.

The uniform and point-mass measures have the intended finite-probability reading, but the current Lean reachability statement returns the lightweight \texttt{Model}, not a \texttt{StrongProbabilityModel} certificate. Its separate strong-promotion obligation is not discharged by the generic convex-path constructor.

The reachability theorem establishes permissiveness across the witness family. The next result identifies the exact local condition under which no categorical change occurs.
